\clearpage
\begin{examplebox}
\begin{center}
\begin{tikzpicture}[font=\small,>=stealth]
\draw[gray!40,rounded corners] (0,0) rectangle (14.2,9);
\node[font=\small\bfseries] at (7.1,8.35) {例1与例2：关于 y 轴对称，截距也跟着变号};

\begin{scope}[shift={(7.1,3.3)},scale=.87]
\draw[->,gray] (-6,0)--(6,0) node[right] {$x$};
\draw[->,gray] (0,-1)--(0,4.6) node[above] {$y$};
\draw[teal!75!black,thick,domain=-.55:.85,samples=70,variable=\t] plot ({3*cosh(\t)},{4*sinh(\t)});
\draw[teal!75!black,thick,domain=-.55:.85,samples=70,variable=\t] plot ({-3*cosh(\t)},{4*sinh(\t)});
\draw[blue!60,dashed] (-5,0)--(3.75,3)--(5,0)--(-3.75,3)--cycle;
\draw[magenta!80!black,very thick] (2.4,0)--(3.75,3);
\draw[magenta!80!black,very thick] (-2.4,0)--(-3.75,3);
\fill (3.75,3) circle(2pt) node[above left] {$P(15/4,3)$};
\fill (-3.75,3) circle(2pt) node[above right] {$Q(-15/4,3)$};
\fill (2.4,0) circle(2pt) node[below left] {$M$};
\fill (-2.4,0) circle(2pt) node[below right] {$N$};
\fill (-5,0) circle(2pt) node[below left] {$F_1$};\fill (5,0) circle(2pt) node[below right] {$F_2$};
\end{scope}
\node at (7.1,.75) {$x_0$ 与 $m$ 同号；左右两支都满足 $x_0m=9$。};

\end{tikzpicture}
\end{center}

\textbf{例1（右支：直接求角平分线交点与切线）。}

双曲线 $\dfrac{x^2}{9}-\dfrac{y^2}{16}=1$，
焦点 $F_1(-5,0),F_2(5,0)$，取点
$P\left(\dfrac{15}{4},3\right)$。
设 $\angle F_1PF_2$ 的内角平分线交 $x$ 轴于 $M(m,0)$。
求 $m$ 及 $PM$ 方程。

\textbf{解：} $a=3,b=4,c=5,x_0=15/4$，直接套用
\[
m=\frac{a^2}{x_0}=\frac{9}{15/4}=\boxed{\frac{12}{5}},
\qquad\boxed{M\left(\frac{12}{5},0\right)}.
\]
对应切线为
\[
\frac{(15/4)x}{9}-\frac{3y}{16}=1
\quad\Longleftrightarrow\quad
\boxed{20x-9y=48}.
\]
\textbf{校验：} $PF_1=37/4,\ PF_2=13/4$，
\[
\frac{F_1M}{MF_2}
=\frac{5+12/5}{5-12/5}=\frac{37}{13}
=\frac{PF_1}{PF_2}.
\]

\medskip
\textbf{例2（左支：完全不用另背一套公式）。}

在同一双曲线上取 $Q\left(-\dfrac{15}{4},3\right)$，
其焦点三角形的顶角内角平分线交轴于 $N(n,0)$。
直接得
\[
\boxed{n=\frac{9}{-15/4}=-\frac{12}{5}},
\qquad\boxed{-20x-9y=48}.
\]
左支的 $n$ 为负，且 $-a<n<0$，与几何位置完全一致。

\medskip
\end{examplebox}
\clearpage
\begin{examplebox}
\begin{center}
\begin{tikzpicture}[font=\small,>=stealth]
\draw[gray!40,rounded corners] (0,0) rectangle (14.2,10);
\node[font=\small\bfseries] at (7.1,9.35) {例3：先固定横坐标，再找上下两个交点};

\begin{scope}[shift={(2.2,5.1)},scale=.48]
\draw[->,gray] (-1,0)--(9,0) node[right] {$x$};
\draw[->,gray] (0,-8)--(0,8) node[above] {$y$};
\draw[teal!75!black,thick,domain=-1.38:1.38,samples=100,variable=\t] plot ({3*cosh(\t)},{4*sinh(\t)});
\draw[dashed,orange!85!black] (6,-7.7)--(6,7.7);
\draw[magenta!80!black,thick] (6,6.9282)--(1.5,0)--(6,-6.9282);
\fill (1.5,0) circle(3pt);\node[below left] at (1.5,0) {$M$};
\fill (6,6.9282) circle(3pt);\draw (6,6.9282)--(7.1,6.9282) node[right] {$P_+$};
\fill (6,-6.9282) circle(3pt);\draw (6,-6.9282)--(7.1,-6.9282) node[right] {$P_-$};
\node[right] at (6,2.4) {$x=6$};
\end{scope}
\node[align=left,anchor=west] at (8.8,6) {$m=3/2$\\[8pt]$x_0=9/m=6$\\[12pt]$P_+=(6,4\sqrt3)$\\[8pt]$P_-=(6,-4\sqrt3)$};
\node at (7.1,.65) {同一个轴交点，通常对应关于 x 轴对称的两个非顶点。};

\end{tikzpicture}
\end{center}

\textbf{例3（反向设点：先知轴截距，再求曲线上点）。}

设 $a=3,b=4$，双曲线上一点 $P(x_0,y_0)$ 的焦点三角形
顶角内角平分线交 $x$ 轴于 $M\left(\dfrac32,0\right)$。
求所有可能的 $P$。

\textbf{解：}从最简结论反推
\[
x_0=\frac{a^2}{m}=\frac{9}{3/2}=6.
\]
代入双曲线
\[
\frac{36}{9}-\frac{y_0^2}{16}=1
\quad\Longrightarrow\quad y_0^2=48.
\]
所以
\[
\boxed{P=(6,4\sqrt3)\quad\text{或}\quad(6,-4\sqrt3)}.
\]
两点都不是顶点，故均满足题设。

\medskip
\end{examplebox}
\clearpage
\begin{examplebox}
\textbf{例4（内切圆切点：分清 $T$ 和 $M$）。}

仍取例1的 $P\left(\frac{15}{4},3\right)$，
作 $\triangle PF_1F_2$ 内切圆，设其与 $x$ 轴切于 $T$。
因为 $P$ 在右支，
\[
\boxed{T=(a,0)=(3,0)},\qquad
\boxed{M=\left(\frac{12}{5},0\right)}.
\]
故 $T\ne M$。进一步，内切圆半径
\[
\rho=\frac{[PF_1F_2]}{s}
=\frac{\frac12\cdot10\cdot3}{\frac12(37/4+13/4+10)}
=\boxed{\frac43},
\]
内心为 $I=(3,4/3)$（点 $P$ 在 $x$ 轴上方）。
这也检验了 $P,I,M$ 三点共线。

\medskip
\textbf{例5（考场变式：切线横截距的范围）。}

设 $P$ 在双曲线的右支且不是顶点，
$PM$ 为顶角内角平分线。因为 $x_0>a$，
\[
0<m=\frac{a^2}{x_0}<a.
\]
所以交点 $M$ 必落在\textbf{原点与右顶点之间}，
并且随着 $P$ 的横坐标变大，$m$ 单调变小。
左支则有 $-a<m<0$。
\end{examplebox}

\clearpage
\begin{examplebox}
\textbf{例6（新增：倒角公式法会出现两个根，哪个才是内角？）。}

双曲线 $\dfrac{x^2}{9}-\dfrac{y^2}{16}=1$，
取 $P\left(\dfrac{15}{4},3\right)$、$F_1(-5,0)$、$F_2(5,0)$。
分别求过 $P$ 的内、外角平分线与 $x$ 轴的交点及其斜率。

\textbf{解：} $a=3,c=5,x_0=\dfrac{15}{4}$，由方法六推导的因式分解式
\[
 (mx_0-a^2)(c^2x_0-a^2m)=0
\]
可得两个候选轴截距
\[
 m_\mathrm{in}=\frac{9}{15/4}=\boxed{\frac{12}{5}},\qquad
 m_\mathrm{out}=\frac{25(15/4)}{9}=\boxed{\frac{125}{12}}.
\]
因为 $-5<\dfrac{12}{5}<5$，而 $\dfrac{125}{12}>5$，
故第一个是\textbf{内角平分线}交点，第二个是\textbf{外角平分线}交点。
相应斜率
\[
 k_\mathrm{in}=\frac{3}{15/4-12/5}=\boxed{\frac{20}{9}},\qquad
 k_\mathrm{out}=\frac{3}{15/4-125/12}=\boxed{-\frac9{20}},
\]
确有 $k_\mathrm{in}k_\mathrm{out}=-1$。

\textbf{检验夹角正切：} $k_1=\dfrac{12}{35},\quad k_2=-\dfrac{12}{5}$，
把 $k=20/9$ 代回式 (6.1)，左右两边均为 $16/15$。
注意：\emph{仅验算正切等式还不够}，必须结合底边内部这一几何条件。
\end{examplebox}
